% ==================================================================
% CHAPTER 3: METHODOLOGY
% ==================================================================
\chapter{Methodology}

\section{Development Model}

The project follows an \textbf{Incremental and Iterative Development Model} with 19 milestones organized into three phases:
\begin{itemize}[leftmargin=1.2cm]
    \item \textbf{Phase I (M1--M5):} Infrastructure setup, authentication, database schema, offline-first CRUD, and duplicate detection.
    \item \textbf{Phase II (M6--M10):} Public dashboard, geo-need engine, feedback module, inventory tracking, and advanced duplicate audit.
    \item \textbf{Phase III (M11--M19):} Full UI refactor, heatmap visualization, pledge coordination system, QA, and documentation.
\end{itemize}

\subsection{Justification}
ReliefMesh required a model that could deliver working increments early and adapt to evolving requirements discovered during prototype testing. A pure waterfall model would have delayed testing until all modules were complete. Scrum was not adopted because the team lacked formal Scrum training and the university course structure imposed fixed deadlines rather than flexible sprint goals.

\subsection{Development Environment}

\begin{table}[H]
\centering
\caption{Development Environment}
\begin{tabular}{@{}ll@{}}
\toprule
\textbf{Component} & \textbf{Technology} \\
\midrule
Frontend & React 18.6.1, Vite, Tailwind CSS 3.4, React Router 6 \\
Backend & Node.js, Express.js 4.18.2 \\
Database & MongoDB 6.0 \\
Authentication & bcrypt.js, jsonwebtoken \\
Offline Storage & IndexedDB (via localforage), Workbox \\
Testing & React Testing Library, Playwright, Jest \\
Deployment & Vercel (Frontend), Render (Backend), MongoDB Atlas (Cloud) \\
\bottomrule
\end{tabular}
\end{table}

\section{Requirements Engineering}

\subsection{Requirement Elicitation}
Requirements were gathered through three channels:
\begin{enumerate}
    \item \textbf{Literature review:} Analysis of prior disaster coordination systems and post-disaster assessment reports for Bangladesh.
    \item \textbf{Hypothetical scenario walkthroughs:} The team modeled disaster response workflows based on published Union Parishad disaster management committee guidelines.
    \item \textbf{System analysis heuristics:} Application of structured analysis techniques (DFD, ERD, use case modeling) to derive functional and non-functional requirements.
\end{enumerate}

\subsection{Functional Requirements}
\begin{enumerate}
    \item \textbf{FR-1:} The system shall allow authorized officials to register affected households with NID, family size, vulnerability flags, ward number, GPS coordinates, and photo.
    \item \textbf{FR-2:} The system shall allow item-level distribution logging with category, quantity, unit, and date.
    \item \textbf{FR-3:} The system shall detect duplicate relief distributions (same household, same category, within 7 days) and flag them with override logging.
    \item \textbf{FR-4:} The system shall compute per-household relief need using Sphere-standard rates with vulnerability multipliers.
    \item \textbf{FR-5:} The system shall provide a public dashboard with aggregated distribution statistics at ward and upazila levels.
    \item \textbf{FR-6:} The system shall display a color-coded heatmap showing served versus remaining-need areas.
    \item \textbf{FR-7:} The system shall support offline data entry with automatic synchronization when connectivity is restored.
    \item \textbf{FR-8:} The system shall support four user roles (UP Official, Upazila Officer, NGO Worker, Public Viewer) with jurisdiction-scoped access.
\end{enumerate}

\subsection{Non-Functional Requirements}
\begin{enumerate}
    \item \textbf{NFR-1:} The system shall function offline for at least 72 hours of continuous use.
    \item \textbf{NFR-2:} The public dashboard shall load within 3 seconds on 3G network.
    \item \textbf{NFR-3:} Duplicate detection shall complete within 2 seconds per distribution.
    \item \textbf{NFR-4:} The system shall support up to 50 concurrent users.
    \item \textbf{NFR-5:} Authentication tokens shall expire after 24 hours.
    \item \textbf{NFR-6:} All relief override actions shall be logged with timestamps and user identity.
\end{enumerate}

\section{System Modeling: Data Flow Diagrams}

Data Flow Diagrams (DFDs) are the central modeling technique used to document the system's functional requirements. The analysis follows a leveled DFD approach, starting from a context diagram and decomposing into detailed diagrams.

\subsection{Context Diagram (Level 0)}

The context diagram shows ReliefMesh as a single process interacting with five external entities:

\begin{figure}[H]
\centering
\begin{tikzpicture}[node distance=2.5cm, auto,
    entity/.style={rectangle, draw, minimum width=2.8cm, minimum height=1cm, align=center},
    process/.style={rectangle, draw, minimum width=2.5cm, minimum height=1.5cm, align=center, rounded corners=2mm},
    dataflow/.style={->, >=stealth, thick}]

    % Central process
    \node[process, fill=blue!10] (system) {\textbf{ReliefMesh}\\System};

    % External entities
    \node[entity, above left=of system] (up) {UP Official};
    \node[entity, below left=of system] (upazila) {Upazila Officer};
    \node[entity, above right=of system] (ngo) {NGO Worker};
    \node[entity, below right=of system] (public) {Public Viewer};
    \node[entity, above=of system, yshift=0.5cm] (admin) {System Admin};

    % Data flows
    \draw[dataflow] (up.east) -- node[above, xshift=-0.3cm] {Register HH} node[below, xshift=-0.3cm] {Log Distrib.} (system.west);
    \draw[dataflow] (system.west) -- node[above, xshift=0.3cm] {Reports} node[below, xshift=0.3cm] {Alerts} (up.west);
    \draw[dataflow] (upazila.east) -- node[above] {View Stats} (system.west);
    \draw[dataflow] (system.west) -- node[below] {Reports} (upazila.west);
    \draw[dataflow] (system.east) -- node[above] {View Need} node[below] {Submit Dist.} (ngo.west);
    \draw[dataflow] (ngo.east) -- node[above] {Alloc. Data} node[below] {} (system.east);
    \draw[dataflow] (public.east) -- node[above] {View Dashboard} (system.east);
    \draw[dataflow] (system.east) -- node[below] {Heatmap} node[below] {} (public.east);

\end{tikzpicture}
\caption{Context Diagram (Level 0 DFD) -- ReliefMesh System}
\label{fig:context_dfd}
\end{figure}

\noindent\textit{Source:} \texttt{diagrams/DFD-context-level-0.drawio}

\subsection{Level 1 DFD}

The Level 1 DFD decomposes the system into four major processes:

\begin{figure}[H]
\centering
\begin{tikzpicture}[node distance=2.8cm,
    proc/.style={rectangle, draw, minimum width=3cm, minimum height=1.5cm, align=center, rounded corners=2mm, fill=blue!5},
    store/.style={rectangle, draw, minimum width=2.5cm, minimum height=1cm, align=center, fill=green!5},
    ext/.style={rectangle, draw, minimum width=2.5cm, minimum height=1cm, align=center, fill=orange!5},
    flow/.style={->, >=stealth, thick}]

    % External entities
    \node[ext] (up) {UP Official};
    \node[ext, below=of up] (upazila) {Upazila Officer};

    % Processes
    \node[proc, right=of up, xshift=1.5cm] (p1) {1.0\\Household Mgmt};
    \node[proc, below=of p1] (p2) {2.0\\Relief Distrib.};
    \node[proc, right=of p1, xshift=2.5cm] (p3) {3.0\\Duplicate Detect};
    \node[proc, below=of p3] (p4) {4.0\\Report \& Dashboard};

    % Data stores
    \node[store, right=of p2, xshift=1.5cm] (d1) {D1\\Households};
    \node[store, right=of p4, xshift=1.5cm] (d2) {D2\\Distributions};

    % Flows: P1
    \draw[flow] (up.east) -- node[above] {HH Data} (p1.west);
    \draw[flow] (p1.east) -- node[above] {HH Record} (d1.west);
    \draw[flow] (p1.south) -- node[left] {HH Info} (p2.north);

    % Flows: P2
    \draw[flow] (p2.east) -- node[above] {Dist. Record} (d2.west);
    \draw[flow] (p2.north) -- node[right, xshift=0.3cm] {Check Dup.} (p3.south);
    \draw[flow] (p3.south) -- node[left, xshift=-0.3cm] {Flag/OK} (p2.north);

    % Flows: P3
    \draw[flow] (p3.west) -- node[above] {Query HH} (d1.east);

    % Flows: P4
    \draw[flow] (p4.north) -- node[right] {Stats} (d2.south);
    \draw[flow] (p4.west) -- node[above] {Dashboard} (upazila.east);

\end{tikzpicture}
\caption{Level 1 Data Flow Diagram}
\label{fig:dfd_l1}
\end{figure}

\noindent\textit{Source:} \texttt{diagrams/DFD-level-1.drawio}

\subsection{Use Case Model}

\begin{figure}[H]
\centering
\begin{tikzpicture}[
    actor/.style={rectangle, draw, minimum width=2.5cm, minimum height=0.8cm, align=center, fill=orange!10},
    uc/.style={ellipse, draw, minimum width=3cm, minimum height=0.8cm, align=center, fill=blue!5},
    boundary/.style={rectangle, draw, dashed, inner sep=0.5cm},
    every node/.style={font=\small}]

    \node[boundary] (boundary) at (1, -2.5) {};
    \node[above] at (boundary.north) {\textbf{ReliefMesh System}};

    % Actors (left side)
    \node[actor] (up) at (-5, 2) {UP Official};
    \node[actor] (upazila) at (-5, 0.5) {Upazila Officer};
    \node[actor] (ngo) at (-5, -1) {NGO Worker};
    \node[actor] (pub) at (-5, -2.5) {Public Viewer};
    \node[actor] (admin) at (-5, -4) {Admin};

    % Use cases (inside boundary)
    \node[uc] (uc1) at (2, 4) {Register Household};
    \node[uc] (uc2) at (4, 2.5) {Log Distribution};
    \node[uc] (uc3) at (2, 1) {Override Duplicate};
    \node[uc] (uc4) at (4, -0.5) {View Dashboard};
    \node[uc] (uc5) at (2, -2) {View Heatmap};
    \node[uc] (uc6) at (4, -3.5) {Submit Feedback};
    \node[uc] (uc7) at (2, -5) {Manage Users};

    % Connections
    \draw (up.east) -- (uc1.west);
    \draw (up.east) -- (uc2.west);
    \draw (up.east) -- (uc3.west);
    \draw (upazila.east) -- (uc4.west);
    \draw (upazila.east) -- (uc5.west);
    \draw (pub.east) -- (uc6.west);
    \draw (admin.east) -- (uc7.west);

\end{tikzpicture}
\caption{Use Case Diagram (simplified view)}
\label{fig:usecase_simple}
\end{figure}

\noindent\textit{Source:} \texttt{diagrams/use-case-diagram.drawio}

\subsection{Entity Relationship Diagram (ERD)}

The logical database schema uses the following entity types:

\begin{figure}[H]
\centering
\begin{tikzpicture}[node distance=1.5cm,
    entity/.style={rectangle, draw, minimum width=3.5cm, minimum height=0.8cm, align=center, fill=yellow!10},
    rel/.style={diamond, draw, align=center, minimum width=2.5cm, minimum height=0.8cm, aspect=2},
    every node/.style={font=\small}]

    % Entities
    \node[entity] (user) {User};
    \node[entity, below=of user] (household) {Household};
    \node[entity, below=of household] (dist) {Distribution};
    \node[entity, right=of household, xshift=2cm] (need) {NeedCalc};
    \node[entity, below=of dist] (pledge) {Pledge};
    \node[entity, right=of dist, xshift=2cm] (feedback) {Feedback};
    \node[entity, right=of user, xshift=2cm] (inventory) {Inventory};

    % Relationships
    \node[rel] (r1) at ($(user)!0.5!(household)$) {registers};
    \node[rel] (r2) at ($(household)!0.5!(dist)$) {receives};
    \node[rel] (r3) at ($(household)!0.5!(need)$) {has};

    % Lines
    \draw (r1) -- (user);
    \draw (r1) -- (household);
    \draw (r2) -- (household);
    \draw (r2) -- (dist);
    \draw (r3) -- (household);
    \draw (r3) -- (need);
    \draw (dist) -- (pledge);
    \draw (dist) -- (feedback);

\end{tikzpicture}
\caption{Entity Relationship Diagram (Logical Schema)}
\label{fig:erd}
\end{figure}

\noindent\textit{Source:} \texttt{diagrams/ER-diagram.drawio}

\section{Algorithm Design: Duplicate Detection}

The duplicate detection algorithm is a core business logic component. It checks whether the same household has received the same category of relief within a configurable lookback window (default 7 days).

\begin{figure}[H]
\centering
\begin{tikzpicture}[node distance=1.2cm, auto,
    block/.style={rectangle, draw, fill=blue!5, minimum width=8cm, minimum height=0.7cm, align=center, font=\small},
    decision/.style={diamond, draw, fill=orange!5, minimum width=3.5cm, minimum height=1.5cm, align=center, font=\small},
    term/.style={rectangle, draw, fill=gray!10, minimum width=5cm, minimum height=0.7cm, rounded corners=3mm, align=center, font=\small},
    arrow/.style={->, >=stealth, thick}]

    % Start
    \node[term] (start) {Input: dist $d(h, c, t)$};
    \node[block, below=of start] (query) {Query distributions for household $h$};
    \node[decision, below=of query, yshift=-0.3cm] (dec1) {Matching $c$?};
    \node[decision, below=of dec1, yshift=-0.3cm] (dec2) {Inside window?};
    \node[block, below=of dec2, yshift=-0.3cm] (log) {Log override};
    \node[term, below=of log, yshift=-0.2cm] (ret) {\color{red} Return DUPLICATE};
    \node[term, right=of dec2, xshift=4cm] (ret2) {\color{green} Return PASS};

    \draw[arrow] (start) -- (query);
    \draw[arrow] (query) -- (dec1);
    \draw[arrow] (dec1) -- node[right] {Yes} (dec2);
    \draw[arrow] (dec1) -| node[left] {No} (ret2);
    \draw[arrow] (dec2) -- node[right] {Yes} (log);
    \draw[arrow] (log) -- (ret);
    \draw[arrow] (dec2) -- node[above] {No} (ret2);
\end{tikzpicture}
\caption{Duplicate Detection Algorithm Flowchart}
\label{alg:dupdetect}
\end{figure}

\section{Architecture}

\subsection{High-Level Architecture}
The system uses a three-tier architecture:
\begin{itemize}[leftmargin=1.2cm]
    \item \textbf{Presentation Tier:} React 18 SPA with Vite bundler, Tailwind CSS, Recharts for visualization, Leaflet for maps.
    \item \textbf{Application Tier:} Express.js REST API with JWT authentication, role-based middleware, and Mongoose ODM.
    \item \textbf{Data Tier:} MongoDB 6.0 with collections for users, households, distributions, pledges, feedback, inventory, and need calculations.
\end{itemize}

\subsection{Offline-First Architecture}
The offline-first layer uses IndexedDB via localforage. Three stores are maintained: \texttt{pending-sync} (queued operations), \texttt{sync-meta} (sync timestamps), and \texttt{offline-data} (cached server data). Background sync triggers when connectivity is restored, sending queued operations with conflict resolution (last-write-wins with manual queue for unresolvable conflicts).

\subsection{Security Architecture}
Authentication uses JWT (24-hour expiry) with bcrypt password hashing. Route-level middleware enforces role-based access: \texttt{upOfficialOnly}, \texttt{upazilaOfficerOnly}, \texttt{ngoWorkerOnly}. Database queries are scoped by jurisdiction. Helmet.js sets security headers and rate limiting prevents abuse.

\section{Testing Strategy}

Three levels of testing:
\begin{enumerate}
    \item \textbf{Unit Tests (Vitest):} 15 test files covering authentication, offline detection, duplicate detection, need calculation, and heatmap logic.
    \item \textbf{E2E Tests (Playwright):} 12 test files covering login, registration, household creation, role access, and visual regression.
    \item \textbf{Manual QA:} 31 structured test cases covering authentication, offline-first data entry, duplicate detection, need calculation, reports, RBAC, and visual regression.
\end{enumerate}

Test results: 10/12 E2E tests passed. 14/15 unit tests passed. 31/31 manual QA tests passed.
